\documentclass[10pt,a4paper]{amsart}
\usepackage[margin=19mm]{geometry}
\usepackage{amsmath,amssymb,mathtools}
\usepackage[scr=rsfs,cal=euler]{mathalfa}
\usepackage{xfrac}
\usepackage[unicode,hidelinks,bookmarks=false]{hyperref}
\newcommand{\ie}{i.e.\ }
\newcommand{\cf}{cf.\ }
\newcommand{\resp}{respectively\ }
\newcommand{\Subsection}{\subsection}
\providecommand{\nobreakdash}{\nobreak\hbox{-}}
\setlength{\emergencystretch}{2em}
\multlinegap0pt
\usepackage{amssymb}
\usepackage[shortlabels]{enumitem}
\usepackage{mathtools}
\usepackage{algorithm}\usepackage{algpseudocode}
\newtheorem{theorem}{Theorem}
\newtheorem{lemma}{Lemma}
\newcommand{\Prob}{\mathbf{P}}
\newcommand{\Z}{\mathbf{Z}}
\newcommand{\bnote}[1]{}
\newcommand{\cardinality}{\#}
\newcommand{\change}[1]{#1}
\newcommand{\lt}{\kappa}
\newcommand{\ncalib}{c}
\newcommand{\param}{\theta}
\newcommand{\paramdist}{\mathbf{Q}}
\newcommand{\predgap}{\hat{g}}
\newcommand{\pst}{\hat{\tau}}
\newcommand{\reals}{\mathbf{R}}
\newcommand{\rollmin}{\bar{g}}
\newcommand{\tol}{\epsilon}
\newcommand{\truegap}{g}
\title{Learning Sequential Mixed Integer Programming Policies}
\author{Stefan Clarke, Bartolomeo Stellato}
\date{Source: arXiv:2602.01476v2 (CC BY 4.0)}
\begin{document}
\maketitle

\noindent\emph{Source and licence.} Learning Sequential Mixed Integer Programming Policies. Version arXiv:2602.01476v2 (CC BY 4.0), distributed by arXiv under the Creative Commons Attribution 4.0 International licence, http://creativecommons.org/licenses/by/4.0/, as recorded in the arXiv licence metadata for this exact version and shown on the abstract page https://arxiv.org/abs/2602.01476v2. Full licence notice: \texttt{licenses/arxiv-2602.01476-CC-BY-4.0.txt}.\par\smallskip

\noindent\textit{Editorial scope.} Counted unit: the article's theorem thm:conformal (source0.tex lines 574--582). The article's complete original proof of it is delivered with it (source0.tex lines 608--647, one \texttt{proof} environment of 40 lines, which the article places away from the statement). No mathematics is written, repaired or re-derived: the statement and the proof are the article's own lines, in the article's own notation, and no retained line is substituted or re-flowed.  Retained uncounted as the source-local context the proof needs: lines 584--591.  Nothing was omitted from any retained span, so the package carries no omission record.  No retained line contains a citation, so the wrapper carries no bibliography and no bibliography entry.  No retained line contains a figure environment, an image-inclusion command or drawing code, so no image is delivered and the package declares no asset dependency.  Editorial formatting only, in addition: an AMS \texttt{amsart} preamble that restates the article's own declarations and macros used by the retained lines, namely the environments \texttt{theorem}, \texttt{lemma} on separate counters, together with the standard packages those lines need.  The retained environments print in this document as Theorem 1, Lemma 1 in order of appearance, whereas the article prints the same items with the numbering of its own full document.  The complete native source of the article as distributed in the arXiv e-print for this exact version is bundled unchanged as \texttt{sources/193791/source0.tex}.
\begin{lemma} \label{lemma:ordering}
    Let $Z_1, \dots, Z_{\ncalib + 1}$ be iid random variables in $\reals$.
    Let $Z_{[1]}, \dots, Z_{[\ncalib]}$ denote $Z_1, \dots, Z_\ncalib$ arranged in increasing order.
    For any $n \in \{1, \dots, \ncalib\}$,
    $$
    \Prob[Z_{\ncalib + 1} \ge Z_{[\ncalib + 1 - n]}] \ge n / (\ncalib + 1).
    $$
\end{lemma}

\begin{theorem} \label{thm:conformal}
    Let $\ncalib \in \Z_+$ and $n \in \{1, \dots, \ncalib\}$.
    Let $\param_1, \dots, \param_\ncalib, \param_{\ncalib+1} \sim \paramdist$ be independent.
    Let
    \begin{equation*}
        \lt = \sup \{ k \ge 0 \mid \#\{i \in \{1, \dots, \ncalib\} \mid \truegap_{\param_i}(\pst_k(\param_i)) \le \tol\} \ge n\}.
    \end{equation*}
    Then $\Prob[\truegap_{\param_{\ncalib+1}}(\pst_{\lt}(\param_{\ncalib+1})) \le \tol] \ge \change{n}/(\ncalib + 1)$.
\end{theorem}

\begin{proof}[Proof of Theorem~\ref{thm:conformal}]
    By definition, $\lt$ is the largest $k \ge 0$ such that at least $n$ of $\{\truegap_{\param_i}(\pst_k(\param_i))\}_{i=1}^\ncalib$ are at most $\tol$.
    Since $\truegap_\param(t)$ is nonincreasing,
    $$
    \truegap_\param(t) \le \tol \iff \truegap_\param^{-1}(\tol) \le t.
    $$
    We also have $\rollmin_\param^{-1} = \predgap_\param^{-1}$, so
    \begin{align}
        \rollmin_\param(t) \ge k \iff \rollmin_\param^{-1}(k) \ge t \iff \predgap_\param^{-1}(k) \ge t. \label{eq:confproof1}
    \end{align}
    % Let 
    % $$
    % m = \#\{i \in \{1, \dots, \ncalib\} \mid \truegap_{\param_i}(\pst_k(\param_i)) \le \tol\}.
    % $$
    Observe that,
    \begin{align*}
    \truegap_{\param_i}^{-1}(\tol) \le \pst_k(\param_i) &\iff \truegap_{\param_i}^{-1}(\tol) \le \rollmin_{\param_i}^{-1}(k) \\&\iff \rollmin_{\param_i}(\truegap_{\param_i}^{-1}(\tol)) \ge k,
    \end{align*}
    where the first equivalence is the definition of $\pst_k$ and the second is by~\eqref{eq:confproof1}.
    % \bnote{do we ever use $N_k$ except for this one time? If not, I would remove the definition}
    % Therefore, if we let 
    % $$
    % N_k = \cardinality\{i \in \{1, \dots, \ncalib\} \mid \rollmin_{\param_i}(\truegap_{\param_i}^{-1}(\tol)) \ge k \},
    % $$
    % by definition,
    Therefore, by definition, $\lt$ is exactly
    \begin{equation*}
        % \lt
        % &= \sup \{ k \ge 0 \mid  \ge n\} \\
        % &= \sup\{k \ge 0 \mid \cardinality\{i \in \{1, \dots, \ncalib\} \mid \truegap_{\param_i}^{-1}(\tol) \le \pst_k(\param_i) \} \ge n \} \\
        % &= \sup\{k \ge 0 \mid \cardinality\{i \in \{1, \dots, \ncalib\} \mid \truegap_{\param_i}^{-1}(\tol) \le \rollmin_{\param_i}^{-1}(k) \} \ge n \} \\
        \sup\{k \ge 0 \mid \cardinality\{i \in \{1, \dots, \ncalib\} \mid \rollmin_{\param_i}(\truegap_{\param_i}^{-1}(\tol)) \ge k \} \ge n \},
    \end{equation*}
    which corresponds to the $n$-th largest element of $\{\rollmin_{\param_i}(\truegap_{\param_i}^{-1}(\tol))\}_{i=1}^\ncalib$.
    By Lemma~\ref{lemma:ordering} applied to $Z_i = \rollmin_{\param_i}(\truegap^{-1}_{\param_i}(\tol))$,
    \begin{align*}
    \Prob[\rollmin_{\param_{\ncalib+1}}(\truegap^{-1}_{\param_{\ncalib + 1}}(\tol)) \ge \lt] &= \Prob[Z_{\ncalib + 1} \ge Z_{[\ncalib + 1 - n]}] \\ & \ge n / (\ncalib + 1).
    \end{align*}
    Reversing the equivalences above yields $\Prob[\truegap_{\param_{\ncalib+1}}(\predgap^{-1}_{\param_{\ncalib+1}}(\lt)) \le \tol] \ge n/(\ncalib + 1)$.
\end{proof}

\end{document}
